% =============================================================================
% ANALISE ESTATISTICA - gRPC x REST (tamanho do payload)
% Autor: Vitor de Freitas Coimbra
% Padrao ABNT identico ao Mapeamento Sistematico (compila rapido no Overleaf Free!)
%
% Compilar no Overleaf: pdflatex -> bibtex -> pdflatex -> pdflatex
% (muito mais leve que biblatex/biber — evita timeout no plano gratuito)
% =============================================================================

\documentclass[a4paper, 12pt]{article}
\usepackage[top=3cm, bottom=2cm, left=3cm, right=2cm]{geometry}
\usepackage[utf8]{inputenc}
\usepackage[T1]{fontenc}
\usepackage{amsmath, amsfonts, amssymb}
\usepackage{graphicx}
\graphicspath{ {./images/} }
\usepackage[brazilian]{babel}
\usepackage{array}
\usepackage{multirow}
\usepackage{enumitem}
\usepackage{float}
\usepackage{caption}
\usepackage{tikz}
\usetikzlibrary{positioning, arrows.meta, shapes.geometric, calc}
\usepackage{tabularx}
\usepackage[table]{xcolor}
\usepackage{titlesec}
\usepackage{changepage}
\usepackage{setspace}
\usepackage[colorlinks=true, allcolors=black]{hyperref}

% Fonte Times New Roman (identica ao mapeamento sistematico)
\usepackage{mathptmx}

% Citacoes ABNT via abntex2cite + bibtex (muito rapido no Overleaf gratuito)
\usepackage[alf]{abntex2cite}
\def\UrlLeft{}\def\UrlRight{}

% Layout identico ao artigo modelo do mapeamento sistematico
\singlespacing
\setlength{\parindent}{1.25cm}
\setlength{\parskip}{0pt}
\pagestyle{empty}

% Secoes e subsecoes
\titleformat{\section}{\fontsize{14}{17}\selectfont\bfseries}{\thesection}{0.5em}{}
\titlespacing*{\section}{0pt}{14pt}{8pt}
\titleformat{\subsection}{\fontsize{13}{16}\selectfont\bfseries}{\thesubsection}{0.5em}{}
\titlespacing*{\subsection}{0pt}{12pt}{6pt}

% Legendas no padrao ABNT: "Tabela 1 -- Texto", acima do elemento, com fonte abaixo
\captionsetup{font=footnotesize, labelfont=bf, labelsep=endash, position=top, skip=4pt}
\DeclareCaptionType{quadro}[Quadro][Lista de quadros]
\newcommand{\fonte}[1]{\par\vspace{3pt}\makebox[\linewidth][l]{\footnotesize Fonte: #1}\par}
\newcommand{\nota}[1]{\par\vspace{1pt}{\footnotesize\noindent Nota: #1\par}}

% Paleta do documento
\definecolor{cabecalho}{RGB}{220,230,242}
\definecolor{destaqueVerde}{RGB}{217,242,217}
\definecolor{bordaVerde}{RGB}{34,120,55}
\definecolor{destaqueAzul}{RGB}{214,230,247}
\definecolor{bordaAzul}{RGB}{30,85,155}
\definecolor{destaqueLaranja}{RGB}{255,224,178}
\definecolor{bordaLaranja}{RGB}{190,95,0}
\definecolor{inativoBg}{RGB}{244,244,244}
\definecolor{inativoFg}{RGB}{145,145,145}
\definecolor{notaBg}{RGB}{255,252,225}
\definecolor{notaBorda}{RGB}{185,150,20}

% Marcador para itens ainda pendentes
\newcommand{\pendente}[1]{\textcolor{red}{[#1]}}

% Macro leve para prints do Minitab: exibe a imagem em ./images/ se existir,
% ou exibe um quadro limpo indicando exatamente qual arquivo subir.
\newcommand{\minitabimg}[3][0.88\linewidth]{%
  \IfFileExists{images/#2}{%
    \includegraphics[width=#1]{#2}%
  }{%
    \fbox{%
      \begin{minipage}[c][2.6cm][c]{0.86\linewidth}
        \centering
        \footnotesize\color{gray!70!black}
        \textbf{[Espa\c{c}o reservado para print do Minitab]}\\[0.35em]
        \texttt{\detokenize{images/#2}}\\[0.25em]
        \scriptsize\textit{#3}
      \end{minipage}%
    }%
  }%
}

% Caixa limpa para saidas textuais literais do Minitab (texto vetorial nitido,
% seguindo o estilo das Figuras 4, 7, 10 e 19 de Sirqueira et al., 2015)
\definecolor{minitabAzul}{RGB}{0,102,164}
\definecolor{minitabBorda}{RGB}{200,208,218}
\definecolor{minitabBg}{RGB}{251,252,254}
\newcommand{\minitabbox}[2][\linewidth]{%
  \fcolorbox{minitabBorda}{minitabBg}{%
    \begin{minipage}{\dimexpr#1-2\fboxsep-2\fboxrule\relax}
      \vspace{2pt}
      #2
      \vspace{2pt}
    \end{minipage}%
  }%
}

\addto\captionsbrazil{\renewcommand{\refname}{REFERÊNCIAS}}

\begin{document}
\citeoption{nbr10520-2023}

% =============================================================================
% TITULO E AUTOR
% =============================================================================
\begin{center}
    {\fontsize{16}{19}\selectfont\bfseries
    ANÁLISE ESTATÍSTICA DO DESEMPENHO DE gRPC E REST EM FUNÇÃO DO TAMANHO DO PAYLOAD EM GO\par}
    \vspace{14pt}
    Vitor de Freitas Coimbra
\end{center}

\vspace{6pt}

% =============================================================================
% RESUMO
% =============================================================================
\begin{adjustwidth}{0.4cm}{0cm}
{\fontsize{13}{16}\selectfont\bfseries RESUMO}
\end{adjustwidth}

\vspace{4pt}

{\small\noindent
\textbf{[Contexto]} A adoção de gRPC em substituição ao estilo arquitetural REST costuma se apoiar na expectativa de menor latência, mas um mapeamento sistemático prévio revelou resultados divergentes conforme o tamanho do payload e escassez de testes estatísticos formais nos estudos empíricos. \textbf{[Objetivo]} Avaliar estatisticamente o desempenho de implementações equivalentes em Go utilizando REST (HTTP/1.1 + JSON) e gRPC (HTTP/2 + \textit{Protocol Buffers}) em quatro tamanhos de payload (1~KB, 100~KB, 1~MB e 4~MB). \textbf{[Método]} Conduziu-se um experimento controlado em macOS (Apple M3) com fase de aquecimento (100 requisições descartadas por cenário), medição com temporizador monotônico de alta resolução e coleta de 1.000 observações por combinação, totalizando 8.000 medições. Os dados foram analisados no software Minitab seguindo integralmente a árvore de decisão de \citeonline{sirqueira2015} ($\alpha = 0{,}05$). \textbf{[Resultados]} A serialização em \textit{Protocol Buffers} reduziu o tamanho trafegado entre 29,08\% e 30,11\%. O teste de Kolmogorov-Smirnov rejeitou a normalidade em todos os oito grupos ($p < 0{,}010$), direcionando a inferência para métodos não paramétricos. O teste de Mann-Whitney identificou diferenças estatisticamente significativas em todos os tamanhos: em 1~KB, ambas as arquiteturas apresentaram empate técnico em magnitude prática (mediana de 0,0405~ms no REST contra 0,0403~ms no gRPC; redução de 0,62\%, $p = 0{,}0115$); a partir de 100~KB o gRPC obteve vantagem expressiva, com reduções medianas de 60,60\% (100~KB), 62,76\% (1~MB) e 62,77\% (4~MB), todos com $p < 0{,}0001$. O teste de Kruskal-Wallis confirmou efeito significativo do tamanho do payload em ambas as arquiteturas ($H = 3.748{,}29$ no REST e $H = 3.749{,}06$ no gRPC; $p < 0{,}0001$). \textbf{[Conclusão]} A vantagem prática do gRPC sobre REST depende fortemente do tamanho do payload: em mensagens muito pequenas (1~KB) sem latência de rede, o custo fixo de enquadramento do HTTP/2/gRPC praticamente neutraliza a economia da serialização binária, enquanto em payloads médios e grandes ($\ge$ 100~KB) o gRPC proporciona ganhos substanciais e estatisticamente consistentes acima de 60\%.\par}

\vspace{10pt}

{\small\noindent\textbf{Palavras-chave:} gRPC; REST; tamanho de payload; análise estatística; Minitab; Engenharia de Software Experimental.\par}

% =============================================================================
% 1. INTRODUCAO
% =============================================================================
\section{INTRODUÇÃO}

A comunicação entre serviços constitui um elemento crítico para o desempenho de sistemas distribuídos baseados em microsserviços. O estilo arquitetural REST (\textit{Representational State Transfer}) \cite{fielding2000}, usualmente implementado sobre HTTP/1.1 com serialização textual JSON, consolidou-se como padrão de mercado pela simplicidade e ampla compatibilidade. Como alternativa voltada a maior eficiência, o gRPC \cite{talwar2016} combina o protocolo de transporte HTTP/2 com a serialização binária do \textit{Protocol Buffers} (Protobuf).

Em etapa anterior desta pesquisa, conduziu-se um mapeamento sistemático da literatura sobre o efeito do tamanho do payload no desempenho de gRPC e REST \cite{coimbra2026}. O mapeamento identificou duas lacunas principais que orientam o presente trabalho: primeiro, os estudos primários não convergem quanto à arquitetura de menor latência em diferentes faixas de tamanho, havendo registros de inversão de vantagem em sentidos opostos \cite{soares2024,kazanavicius2023,bolanowski2022}; segundo, apenas um dos dez estudos analisados aplicou testes estatísticos formais de hipótese, enquanto os demais restringiram-se a estatísticas descritivas.

Na Engenharia de Software Experimental, conclusões empíricas confiáveis dependem de coleta controlada, volume amostral adequado e escolha fundamentada dos testes estatísticos \cite{wohlin2012,sirqueira2015}. Diante disso, este trabalho apresenta um estudo experimental em linguagem Go comparando REST (HTTP/1.1 + JSON) e gRPC (HTTP/2 + Protobuf) em quatro faixas de tamanho de payload (1~KB, 100~KB, 1~MB e 4~MB), aplicando integralmente a árvore de decisão estatística proposta por \citeonline{sirqueira2015} por meio do software Minitab sobre 8.000 observações coletadas.

O restante do documento está organizado da seguinte forma: a Seção 2 detalha o planejamento experimental, a instrumentação e o protocolo estatístico adotado; a Seção 3 apresenta os resultados descritivos, a verificação de pressupostos e os testes de hipótese no Minitab; a Seção 4 discute os achados frente ao mapeamento sistemático e aborda as ameaças à validade; e a Seção 5 apresenta a conclusão e os próximos passos da pesquisa.

% =============================================================================
% 2. METODOLOGIA EXPERIMENTAL
% =============================================================================
\section{METODOLOGIA EXPERIMENTAL}

O experimento foi conduzido conforme o processo de definição, planejamento, execução e análise estabelecido por \citeonline{wohlin2012} e \citeonline{sirqueira2015}.

\subsection{Objetivo, Questões e Hipóteses}

O estudo tem como objetivo analisar as arquiteturas REST (HTTP/1.1) e gRPC (HTTP/2) com o propósito de avaliar estatisticamente sua latência de requisição-resposta em função do tamanho do payload, no contexto de serviços implementados em Go. Para isso, definiram-se duas questões experimentais avaliadas ao nível de significância $\alpha = 0{,}05$:

\begin{itemize}[noitemsep, topsep=4pt]
  \item \textbf{QP1 (1 fator, 2 tratamentos — Arquitetura por Payload):} Há diferença estatisticamente significativa na latência entre REST e gRPC em cada tamanho de payload (1~KB, 100~KB, 1~MB e 4~MB)?
  \begin{itemize}[noitemsep]
    \item $H_{0,\text{arq}}$: Não há diferença significativa entre as latências de REST e gRPC ($\eta_{\text{REST}} = \eta_{\text{gRPC}}$).
    \item $H_{1,\text{arq}}$: Há diferença significativa entre as latências de REST e gRPC ($\eta_{\text{REST}} \neq \eta_{\text{gRPC}}$).
  \end{itemize}
  \item \textbf{QP2 (1 fator, 4 tratamentos — Tamanho de Payload por Arquitetura):} A variação do tamanho do payload (1~KB, 100~KB, 1~MB e 4~MB) produz diferença estatisticamente significativa na latência dentro de cada arquitetura?
  \begin{itemize}[noitemsep]
    \item $H_{0,\text{pay}}$: Não há diferença significativa entre as latências dos quatro tamanhos de payload ($\eta_{\text{1KB}} = \eta_{\text{100KB}} = \eta_{\text{1MB}} = \eta_{\text{4MB}}$).
    \item $H_{1,\text{pay}}$: Pelo menos um tamanho de payload apresenta latência estatisticamente distinta dos demais.
  \end{itemize}
\end{itemize}

\subsection{Variáveis, Instrumentação e Ambiente de Execução}

As variáveis independentes (fatores) são a \textbf{Arquitetura} (escala nominal, 2 tratamentos: REST e gRPC) e o \textbf{Tamanho do Payload} (escala ordinal/razão, 4 tratamentos: 1~KB, 100~KB, 1~MB e 4~MB). A variável dependente principal é a \textbf{latência fim-a-fim da requisição} (\texttt{Latencia\_ms}), medida em milissegundos (escala de razão), acompanhada pelo tamanho serializado em bytes (\texttt{wire\_bytes}) como métrica auxiliar.

Foram implementados dois programas independentes em Go, um para REST (\texttt{cmd/rest\_bench}) e outro para gRPC (\texttt{cmd/grpc\_bench}), compartilhando um gerador determinístico (\texttt{internal/dataset}) de objetos estruturados de domínio (\texttt{Item}, com oito campos tipados: inteiro de 64 bits, quatro cadeias textuais, ponto flutuante, inteiro de 32 bits e booleano). A quantidade de itens foi calibrada para atingir aproximadamente 1~KB (4 itens), 100~KB (378 itens), 1~MB (3.845 itens) e 4~MB (15.348 itens) em JSON, garantindo equivalência exata de conteúdo lógico entre REST e gRPC.

Para controle metodológico, ambas as arquiteturas operaram sobre conexão TCP persistente em \textit{loopback} local (\texttt{127.0.0.1}), sem criptografia TLS e sem compressão gzip. No REST, forçou-se HTTP/1.1 (\texttt{ForceAttemptHTTP2: false}) com \textit{keep-alive}; no gRPC (HTTP/2), elevou-se o limite máximo de mensagem para 32~MB para acomodar o payload de 4~MB. Antes de cada cenário, executou-se coleta de lixo explícita (\texttt{runtime.GC()}) e aquecimento (\textit{warm-up}) de 100 requisições descartadas, seguido de \textbf{1.000 medições válidas por cenário}, perfazendo \textbf{8.000 observações totais}. A medição temporal utilizou temporizador monotônico de alta resolução (\texttt{mach\_absolute\_time} no macOS; \texttt{QueryPerformanceCounter} no Windows), evitando erros de quantização em requisições sub-milissegundo. O Quadro~1 resume o ambiente utilizado na coleta experimental.

\begin{table}[H]
\centering
\captionsetup{type=quadro}
\caption{Especificações do ambiente experimental}
\small
\renewcommand{\arraystretch}{1.35}
\begin{tabularx}{\textwidth}{|>{\raggedright\arraybackslash}p{4.3cm}|>{\raggedright\arraybackslash}X|}
\hline
\rowcolor{cabecalho}
\multicolumn{1}{|c|}{\textbf{Componente}} & \multicolumn{1}{c|}{\textbf{Especificação}} \\
\hline
Sistema Operacional & macOS 15.6.1 (Build 24G90) \\
\hline
Processador & Apple M3 (8 núcleos físicos / 8 threads) \\
\hline
Memória RAM & 16,00 GB \\
\hline
Linguagem e Bibliotecas & Go 1.27.4 (\texttt{darwin/arm64}); \texttt{google.golang.org/grpc v1.71.0}; \texttt{protobuf v1.36.5}; \texttt{protoc 36.2} \\
\hline
Rede e Transporte & Loopback TCP local (\texttt{127.0.0.1}), sem TLS, conexão persistente \\
\hline
Volume de Dados & 100 requisições de aquecimento + 1.000 medições por cenário (8.000 amostras) \\
\hline
Análise Estatística & Minitab ($\alpha = 0{,}05$) \\
\hline
\end{tabularx}
\fonte{elaborado pelo autor (2026).}
\end{table}

\subsection{Protocolo Estatístico e Árvore de Decisão Aplicada}

A análise no Minitab seguiu rigorosamente a árvore de decisão proposta por \citeonline{sirqueira2015}. A Figura~1 adapta o fluxograma original destacando em cores os caminhos efetivamente percorridos neste experimento (e mantendo esmaecidos em cinza os ramos descartados), explicitando em cada nó de decisão a justificativa estatística e os grupos avaliados.

\begin{figure}[H]
\centering
\caption{Árvore de decisão estatística de \citeonline{sirqueira2015} adaptada com os caminhos percorridos e as justificativas deste experimento}
\resizebox{\textwidth}{!}{%
\begin{tikzpicture}[
  font=\scriptsize,
  startend/.style={ellipse, draw=bordaLaranja, fill=destaqueLaranja, thick, minimum width=1.8cm, minimum height=0.65cm, align=center, font=\scriptsize\bfseries},
  decAct/.style={diamond, aspect=2.1, draw=bordaVerde, fill=destaqueVerde, thick, text width=2.45cm, align=center, inner sep=1pt, font=\scriptsize\bfseries},
  decInact/.style={diamond, aspect=2.1, draw=inativoFg, fill=inativoBg, dashed, text=inativoFg, text width=2.35cm, align=center, inner sep=1pt},
  procAct/.style={rectangle, rounded corners=2pt, draw=bordaAzul, fill=destaqueAzul, thick, text width=3.45cm, align=center, inner sep=4pt, font=\scriptsize\bfseries},
  procInact/.style={rectangle, rounded corners=2pt, draw=inativoFg, fill=inativoBg, dashed, text=inativoFg, text width=3.25cm, align=center, inner sep=4pt},
  testAct/.style={rectangle, rounded corners=2pt, draw=bordaLaranja, fill=destaqueLaranja, very thick, text width=3.5cm, align=center, inner sep=4pt, font=\scriptsize\bfseries},
  testInact/.style={rectangle, rounded corners=2pt, draw=inativoFg, fill=inativoBg, dashed, text=inativoFg, text width=2.25cm, align=center, inner sep=3.5pt},
  justif/.style={rectangle, rounded corners=2pt, draw=notaBorda, fill=notaBg, text width=4.0cm, align=left, inner sep=3.5pt, font=\tiny},
  arrAct/.style={-{Stealth[length=2.3mm]}, very thick, draw=bordaAzul},
  arrInact/.style={-{Stealth[length=1.8mm]}, dashed, draw=inativoFg}
]

% Linha 1: Inicio -> N > 30? -> K-S
\node[startend] (ini) {Início};
\node[decAct, right=0.8cm of ini] (n30) {A amostra possui\\mais de 30 valores?};
\node[procAct, right=1.8cm of n30] (ks) {Teste de Kolmogorov-Smirnov (K-S)\\para verificar normalidade};
\node[justif, right=0.5cm of ks] (j1) {\textbf{Decisão (Sim):} Cada grupo possui $N = 1.000$ amostras ($8.000$ no total), logo $N > 30$. Descarta-se Shapiro-Wilk.};

% Linha 2: Shapiro-Wilk (inativo) e Normal para todos? (ativo)
\node[procInact, below=0.85cm of n30] (sw) {Teste de Shapiro-Wilk\\(não utilizado: $N > 30$)};
\node[decAct, below=0.85cm of ks] (norm) {A distribuição é\\normal para todos\\os valores?};
\node[justif, right=0.5cm of norm] (j2) {\textbf{Decisão (Não):} K-S rejeitou normalidade em todos os 8 grupos ($p < 0{,}010 < \alpha = 0{,}05$). Segue direto ao ramo não paramétrico.};

% Linha 3: Nao parametrico (ativo) vs Levene (inativo / complementar)
\node[procAct, below=1.45cm of sw] (nonparam) {Utilizar teste\\não paramétrico};
\node[procInact, below=0.95cm of norm] (levene) {Teste de Levene\\(homocedasticidade)};
\node[decInact, below=0.65cm of levene] (homo) {A amostra é\\homocedástica?};
\node[justif, right=0.5cm of levene, text width=4.0cm] (j3) {\textbf{Verificação complementar:} Embora dispensado após falha na normalidade, Levene também apontou heterocedasticidade em 1~KB, 100~KB e 1~MB ($p < 0{,}05$).};

% Linha 4: Ramo parametrico (inativo) vs Fatores/Tratamentos (ativo)
\node[decAct, below=1.05cm of nonparam] (treatNP) {Quantos fatores\\e tratamentos o\\estudo possui?};
\node[procInact, below=0.65cm of homo] (param) {Utilizar teste\\paramétrico};
\node[decInact, below=0.65cm of param] (treatP) {Quantos fatores\\e tratamentos?};

% Linha 5: Testes finais
\node[testAct, below left=1.25cm and -0.8cm of treatNP] (mw) {Teste de Mann-Whitney\\[2pt]
  \normalfont\tiny\textbf{Análise 1 (1 fator, 2 tratamentos):}\\
  \tiny Fator: \textit{Arquitetura} (grupos: \textbf{REST} vs. \textbf{gRPC}),\\
  \tiny aplicado em 1~KB, 100~KB, 1~MB e 4~MB.};

\node[testAct, below right=1.25cm and -0.8cm of treatNP] (kw) {Teste de Kruskal-Wallis\\(+ Mann-Whitney par-a-par)\\[2pt]
  \normalfont\tiny\textbf{Análise 2 (1 fator, 4 tratamentos):}\\
  \tiny Fator: \textit{Payload} (grupos: \textbf{1~KB, 100~KB, 1~MB, 4~MB}),\\
  \tiny avaliado dentro de REST e de gRPC.};

\node[testInact, below left=0.85cm and -0.35cm of treatP] (ttest) {Teste T\\(não aplicável)};
\node[testInact, below right=0.85cm and -0.35cm of treatP] (anova) {ANOVA + Tukey\\(não aplicável)};

% Termino centralizado abaixo de mw e kw (sem encostar no ramo inativo)
\node[startend] (fim) at ($(mw.south)!0.5!(kw.south) + (0,-1.05cm)$) {Término};

% Setas ativas (sem cruzar nenhum bloco!)
\draw[arrAct] (ini) -- (n30);
\draw[arrAct] (n30) -- node[above, font=\scriptsize\bfseries, text=bordaVerde] {Sim ($N=1000$)} (ks);
\draw[arrAct] (ks) -- (norm);
\draw[arrAct] (norm.south west) -- ++(-0.55cm,-0.35cm) coordinate (midnp) -| (nonparam.north)
  node[pos=0.22, above, font=\scriptsize\bfseries, text=bordaVerde] {Não ($p < 0{,}010$)};
\draw[arrAct] (nonparam) -- (treatNP);
\draw[arrAct] (treatNP.south west) -- node[left, pos=0.45, font=\tiny\bfseries, text=bordaAzul, align=right] {1 fator,\\2 tratamentos} (mw.north);
\draw[arrAct] (treatNP.south east) -- node[right, pos=0.45, font=\tiny\bfseries, text=bordaAzul, align=left] {1 fator,\\4 tratamentos ($>2$)} (kw.north);
\draw[arrAct] (mw.south) |- (fim.west);
\draw[arrAct] (kw.south) |- (fim.east);

% Setas inativas (caminhos descartados do fluxograma original)
\draw[arrInact] (n30) -- node[left, font=\tiny, text=inativoFg] {Não} (sw);
\draw[arrInact] (sw.east) -- (norm.west);
\draw[arrInact] (norm) -- node[right, font=\tiny, text=inativoFg] {Sim} (levene);
\draw[arrInact] (levene) -- (homo);
\draw[arrInact] (homo.west) -| node[pos=0.25, above, font=\tiny, text=inativoFg] {Não} (nonparam.east);
\draw[arrInact] (homo) -- node[right, font=\tiny, text=inativoFg] {Sim} (param);
\draw[arrInact] (param) -- (treatP);
\draw[arrInact] (treatP.south west) -- node[left, font=\tiny, text=inativoFg] {2} (ttest.north);
\draw[arrInact] (treatP.south east) -- node[right, font=\tiny, text=inativoFg] {$>2$} (anova.north);

% Linhas tracejadas conectando caixas de justificativa
\draw[densely dotted, draw=notaBorda, thick] (ks.east) -- (j1.west);
\draw[densely dotted, draw=notaBorda, thick] (norm.east) -- (j2.west);
\draw[densely dotted, draw=notaBorda, thick] (levene.east) -- (j3.west);

\end{tikzpicture}%
}
\fonte{adaptado de \citeonline{sirqueira2015} pelo autor (2026).}
\nota{Blocos coloridos com setas contínuas representam o caminho seguido neste estudo; blocos cinzas tracejados representam os ramos originais da árvore descartados após a verificação dos pressupostos.}
\end{figure}

% =============================================================================
% 3. RESULTADOS E ANALISE ESTATISTICA
% =============================================================================
\section{RESULTADOS E ANÁLISE ESTATÍSTICA}

\subsection{Tamanho das Mensagens Serializadas (JSON vs. Protocol Buffers)}

A Tabela~1 apresenta o tamanho real em bytes trafegado para o mesmo conjunto lógico de registros em cada faixa de payload. A serialização binária em \textit{Protocol Buffers} reduziu o volume de dados entre \textbf{29,08\% e 30,11\%} em todas as faixas avaliadas.

\begin{table}[H]
\centering
\caption{Comparação do tamanho serializado das mensagens na rede (JSON vs. \textit{Protocol Buffers})}
\small
\renewcommand{\arraystretch}{1.35}
\begin{tabularx}{\textwidth}{>{\raggedright\arraybackslash}X*{4}{>{\raggedleft\arraybackslash}p{2.9cm}}}
\hline
\rowcolor{cabecalho}
\multicolumn{1}{c}{\textbf{Payload}} & \multicolumn{1}{c}{\textbf{Qtd. Itens}} & \multicolumn{1}{c}{\textbf{REST / JSON (B)}} & \multicolumn{1}{c}{\textbf{gRPC / Protobuf (B)}} & \multicolumn{1}{c}{\textbf{Redução (\%)}} \\
\hline
1 KB   & 4      & 1.106     & 773       & 30,11\% \\
100 KB & 378    & 102.560   & 72.738    & 29,08\% \\
1 MB   & 3.845  & 1.048.750 & 740.986   & 29,35\% \\
4 MB   & 15.348 & 4.195.335 & 2.958.144 & 29,49\% \\
\hline
\end{tabularx}
\fonte{elaborado pelo autor (2026).}
\end{table}

\subsection{Estatística Descritiva e Análise Exploratória (Boxplots)}

As estatísticas descritivas geradas no Minitab (\textit{Stat $\rightarrow$ Basic Statistics $\rightarrow$ Display Descriptive Statistics}) para os oito grupos experimentais ($N = 1.000$ por grupo) estão consolidadas na Tabela~2, seguindo o formato tabular adotado por \citeonline{sirqueira2015}. Os \textit{boxplots} comparativos por tamanho de payload (\textit{Graph $\rightarrow$ Boxplot $\rightarrow$ One Y / With Groups}) são apresentados na Figura~2.

\begin{table}[H]
\centering
\caption{Estatísticas descritivas da latência (ms) por payload e arquitetura ($N = 1.000$ por grupo)}
\footnotesize
\setlength{\tabcolsep}{3.5pt}
\renewcommand{\arraystretch}{1.3}
\begin{tabularx}{\textwidth}{ll*{9}{>{\raggedleft\arraybackslash}X}}
\hline
\rowcolor{cabecalho}
\textbf{Payload} & \textbf{Arq.} & \textbf{Média} & \textbf{Mediana} & \textbf{DP} & \textbf{Var.} & \textbf{Mín.} & \textbf{Q1} & \textbf{Q3} & \textbf{Máx.} & \textbf{IQR} \\
\hline
\multirow{2}{*}{1 KB}
  & REST & 0,0457 & 0,0405 & 0,0139 & 0,0002 & 0,0302 & 0,0378 & 0,0530 & 0,2392 & 0,0151 \\
  & gRPC & 0,0453 & 0,0403 & 0,0132 & 0,0002 & 0,0285 & 0,0363 & 0,0540 & 0,1825 & 0,0177 \\
\hline
\multirow{2}{*}{100 KB}
  & REST & 0,5108 & 0,4971 & 0,0644 & 0,0041 & 0,4303 & 0,4669 & 0,5255 & 0,8532 & 0,0586 \\
  & gRPC & 0,2035 & 0,1958 & 0,0446 & 0,0020 & 0,1854 & 0,1927 & 0,2000 & 0,7885 & 0,0073 \\
\hline
\multirow{2}{*}{1 MB}
  & REST & 4,7038 & 4,6778 & 0,5308 & 0,2818 & 4,3140 & 4,5392 & 4,8066 & 18,8021 & 0,2674 \\
  & gRPC & 1,7457 & 1,7422 & 0,1419 & 0,0201 & 1,3694 & 1,7167 & 1,7753 & 2,3570  & 0,0587 \\
\hline
\multirow{2}{*}{4 MB}
  & REST & 18,7723 & 18,7834 & 1,2466 & 1,5539 & 17,2046 & 18,6120 & 18,9868 & 47,0041 & 0,3748 \\
  & gRPC & 7,0154  & 6,9928  & 0,6973 & 0,4862 & 5,9417  & 6,8828  & 7,2771  & 25,6896 & 0,3943 \\
\hline
\end{tabularx}
\fonte{elaborado pelo autor (2026), com base nos dados processados no Minitab (\texttt{GRPCxREST.mpx}).}
\end{table}

\begin{figure}[H]
\centering
\caption{Boxplots da latência (ms) comparando REST e gRPC em cada tamanho de payload}
\begin{minipage}[t]{0.48\textwidth}
  \centering
  \minitabimg[\linewidth]{boxplot_1kb.png}{Boxplot REST vs. gRPC (1 KB)}
  \par\vspace{2pt}{\footnotesize (a) Payload de 1 KB}
\end{minipage}\hfill
\begin{minipage}[t]{0.48\textwidth}
  \centering
  \minitabimg[\linewidth]{boxplot_100kb.png}{Boxplot REST vs. gRPC (100 KB)}
  \par\vspace{2pt}{\footnotesize (b) Payload de 100 KB}
\end{minipage}
\vspace{8pt}
\begin{minipage}[t]{0.48\textwidth}
  \centering
  \minitabimg[\linewidth]{boxplot_1mb.png}{Boxplot REST vs. gRPC (1 MB)}
  \par\vspace{2pt}{\footnotesize (c) Payload de 1 MB}
\end{minipage}\hfill
\begin{minipage}[t]{0.48\textwidth}
  \centering
  \minitabimg[\linewidth]{boxplot_4mb.png}{Boxplot REST vs. gRPC (4 MB)}
  \par\vspace{2pt}{\footnotesize (d) Payload de 4 MB}
\end{minipage}
\fonte{elaborado pelo autor (2026) no Minitab.}
\end{figure}

Na maioria das distribuições observa-se assimetria positiva (média superior à mediana) decorrente da cauda superior de latência. Nos \textit{boxplots}, nota-se ampla sobreposição em 1~KB com medianas praticamente idênticas (0,0405~ms no REST contra 0,0403~ms no gRPC), seguida de clara separação favorável ao gRPC em 100~KB, 1~MB e 4~MB.

\subsection{Análise de Outliers}

Aplicou-se o teste de Grubbs no Minitab (\textit{Stat $\rightarrow$ Basic Statistics $\rightarrow$ Outlier Test}, $\alpha = 0{,}05$), complementado pela contagem pelos critérios interquartis ($Q_3 + 1{,}5\cdot\text{IQR}$ e $Q_3 + 3\cdot\text{IQR}$), conforme resumido na Tabela~3 e reproduzido na saída textual do Minitab da Figura~3.

\begin{table}[H]
\centering
\caption{Análise de \textit{outliers} pelo teste de Grubbs e pelo critério interquartil (IQR)}
\small
\renewcommand{\arraystretch}{1.3}
\begin{tabularx}{\textwidth}{>{\raggedright\arraybackslash}p{1.3cm}>{\raggedright\arraybackslash}p{1.1cm}*{5}{>{\raggedleft\arraybackslash}X}}
\hline
\rowcolor{cabecalho}
\textbf{Payload} & \textbf{Arq.} & \textbf{Grubbs $G$} & \textbf{$p$-valor} & \textbf{Extremo (ms)} & \textbf{$1{,}5\cdot\text{IQR}$} & \textbf{$3{,}0\cdot\text{IQR}$} \\
\hline
\multirow{2}{*}{1 KB}   & REST & 13,89 & $< 0{,}001$ & 0,2392  & 17 (1,7\%)   & 5 (0,5\%)    \\
                        & gRPC & 10,42 & $< 0{,}001$ & 0,1825  & 11 (1,1\%)   & 4 (0,4\%)    \\
\hline
\multirow{2}{*}{100 KB} & REST & 5,32  & $< 0{,}001$ & 0,8532  & 97 (9,7\%)   & 20 (2,0\%)   \\
                        & gRPC & 13,10 & $< 0{,}001$ & 0,7885  & 78 (7,8\%)   & 59 (5,9\%)   \\
\hline
\multirow{2}{*}{1 MB}   & REST & 26,56 & $< 0{,}001$ & 18,8021 & 5 (0,5\%)    & 3 (0,3\%)    \\
                        & gRPC & 4,31  & 0,015       & 2,3570  & 251 (25,1\%) & 154 (15,4\%) \\
\hline
\multirow{2}{*}{4 MB}   & REST & 22,65 & $< 0{,}001$ & 47,0041 & 86 (8,6\%)   & 36 (3,6\%)   \\
                        & gRPC & 26,78 & $< 0{,}001$ & 25,6896 & 74 (7,4\%)   & 1 (0,1\%)    \\
\hline
\end{tabularx}
\fonte{elaborado pelo autor (2026).}
\end{table}

\begin{figure}[H]
\centering
\caption{Saída do teste de \textit{outliers} (Grubbs) obtida no Minitab ($\alpha = 0{,}05$)}
\minitabbox[0.95\textwidth]{%
  \ttfamily\footnotesize
  \textbf{\color{minitabAzul}Outlier Test: REST\_1KB, gRPC\_1KB, REST\_100KB, gRPC\_100KB, REST\_1MB, gRPC\_1MB, REST\_4MB, gRPC\_4MB}\\[4pt]
  \textbf{Method}\\
  Null hypothesis \hspace{1.8cm} All data values come from the same normal population\\
  Alternative hypothesis \hspace{0.55cm} Smallest or largest data value is an outlier\\
  Significance level \hspace{1.3cm} $\alpha = 0.05$\\[5pt]
  \textbf{Grubbs' Test}\\[2pt]
  \begin{tabular}{@{}lrrrrrrr@{}}
  \textbf{Variable} & \textbf{N} & \textbf{Mean} & \textbf{StDev} & \textbf{Min} & \textbf{Max} & \textbf{G} & \textbf{P} \\
  \hline
  REST\_1KB   & 1000 & 0.045697 & 0.013926 & 0.030166 & 0.239166 & 13.89 & 0.000 \\
  gRPC\_1KB   & 1000 & 0.045292 & 0.013165 & 0.028541 & 0.182500 & 10.42 & 0.000 \\
  REST\_100KB & 1000 & 0.51075  & 0.06440  & 0.43033  & 0.85321  &  5.32 & 0.000 \\
  gRPC\_100KB & 1000 & 0.20345  & 0.04465  & 0.18542  & 0.78850  & 13.10 & 0.000 \\
  REST\_1MB   & 1000 & 4.7038   & 0.5308   & 4.3140   & 18.8021  & 26.56 & 0.000 \\
  gRPC\_1MB   & 1000 & 1.7457   & 0.1419   & 1.3694   & 2.3570   &  4.31 & 0.015 \\
  REST\_4MB   & 1000 & 18.772   & 1.247    & 17.205   & 47.004   & 22.65 & 0.000 \\
  gRPC\_4MB   & 1000 & 7.0154   & 0.6973   & 5.9417   & 25.6896  & 26.78 & 0.000 \\
  \end{tabular}
}
\fonte{elaborado pelo autor (2026), transcrito da saída do Minitab (\texttt{GRPCxREST.mpx}).}
\end{figure}

De acordo com \citeonline{sirqueira2015}, \textit{outliers} decorrentes de valores fisicamente possíveis da população não devem ser descartados arbitrariamente. Como os picos observados correspondem à latência de cauda (\textit{tail latency}) causada por pausas do coletor de lixo do Go e escalonamento do sistema operacional, optou-se por mantê-los na amostra e conduzir a inferência por métodos não paramétricos baseados em postos e medianas, naturalmente robustos a valores extremos.

\subsection{Verificação dos Pressupostos: Normalidade e Homocedasticidade}

Dado $N = 1.000 > 30$ em todos os grupos, aplicou-se o \textbf{Teste de Kolmogorov-Smirnov (K-S)} (\textit{Stat $\rightarrow$ Basic Statistics $\rightarrow$ Normality Test}) para testar $H_0$: a amostra segue distribuição normal. Complementarmente, verificou-se a igualdade de variâncias pelo \textbf{Teste de Levene} (\textit{Stat $\rightarrow$ Basic Statistics $\rightarrow$ 2 Variances}). A Tabela~4 sintetiza os resultados, ilustrados nas Figuras~4 e 5.

\begin{table}[H]
\centering
\caption{Resultados dos testes de normalidade (Kolmogorov-Smirnov) e homocedasticidade (Levene) ($\alpha = 0{,}05$)}
\small
\renewcommand{\arraystretch}{1.3}
\begin{tabularx}{\textwidth}{>{\raggedright\arraybackslash}p{1.5cm}>{\raggedright\arraybackslash}p{1.2cm}*{5}{>{\centering\arraybackslash}X}}
\hline
\rowcolor{cabecalho}
\textbf{Payload} & \textbf{Arq.} & \textbf{K-S ($D$)} & \textbf{K-S $p$-valor} & \textbf{Normal?} & \textbf{Levene $p$-valor} & \textbf{Homoced.?} \\
\hline
\multirow{2}{*}{1 KB}   & REST & 0,209 & $< 0{,}010$ & Não & \multirow{2}{*}{0,0433}      & \multirow{2}{*}{Não} \\
                        & gRPC & 0,182 & $< 0{,}010$ & Não & & \\
\hline
\multirow{2}{*}{100 KB} & REST & 0,175 & $< 0{,}010$ & Não & \multirow{2}{*}{$< 0{,}001$} & \multirow{2}{*}{Não} \\
                        & gRPC & 0,372 & $< 0{,}010$ & Não & & \\
\hline
\multirow{2}{*}{1 MB}   & REST & 0,262 & $< 0{,}010$ & Não & \multirow{2}{*}{$< 0{,}001$} & \multirow{2}{*}{Não} \\
                        & gRPC & 0,198 & $< 0{,}010$ & Não & & \\
\hline
\multirow{2}{*}{4 MB}   & REST & 0,319 & $< 0{,}010$ & Não & \multirow{2}{*}{0,2690}      & \multirow{2}{*}{Sim} \\
                        & gRPC & 0,212 & $< 0{,}010$ & Não & & \\
\hline
\end{tabularx}
\fonte{elaborado pelo autor (2026).}
\end{table}

\begin{figure}[H]
\centering
\caption{Gráficos de probabilidade normal com teste de Kolmogorov-Smirnov no Minitab (exemplos em 1 KB e 1 MB)}
\begin{minipage}[t]{0.48\textwidth}
  \centering
  \minitabimg[\linewidth]{ks_rest_1kb.png}{Teste K-S: REST 1 KB}
  \par\vspace{2pt}{\footnotesize (a) REST — 1 KB}
\end{minipage}\hfill
\begin{minipage}[t]{0.48\textwidth}
  \centering
  \minitabimg[\linewidth]{ks_grpc_1kb.png}{Teste K-S: gRPC 1 KB}
  \par\vspace{2pt}{\footnotesize (b) gRPC — 1 KB}
\end{minipage}
\vspace{8pt}
\begin{minipage}[t]{0.48\textwidth}
  \centering
  \minitabimg[\linewidth]{ks_rest_1mb.png}{Teste K-S: REST 1 MB}
  \par\vspace{2pt}{\footnotesize (c) REST — 1 MB}
\end{minipage}\hfill
\begin{minipage}[t]{0.48\textwidth}
  \centering
  \minitabimg[\linewidth]{ks_grpc_1mb.png}{Teste K-S: gRPC 1 MB}
  \par\vspace{2pt}{\footnotesize (d) gRPC — 1 MB}
\end{minipage}
\fonte{elaborado pelo autor (2026) no Minitab.}
\end{figure}

\begin{figure}[H]
\centering
\caption{Verificação complementar de homocedasticidade pelo teste de Levene no Minitab (payload de 1 MB)}
\minitabimg[0.78\linewidth]{levene_1mb.png}{Stat -> Basic Statistics -> 2 Variances (1 MB)}
\fonte{elaborado pelo autor (2026) no Minitab.}
\end{figure}

Como todos os grupos violam a normalidade ($p < 0{,}010$) — além de três dos quatro tamanhos violarem também a homocedasticidade ($p < 0{,}05$) —, a árvore de decisão (Figura~1) determina o emprego exclusivo de \textbf{testes não paramétricos}. Ressalta-se que a ausência de normalidade é o comportamento esperado em benchmarks de latência de software: tempos de resposta possuem um limite inferior físico rígido (determinado pelo custo mínimo de CPU, serialização e pilha TCP), mas apresentam cauda longa à direita causada por preempção do sistema operacional, pausas do coletor de lixo e contenção de \textit{cache}. Com amostras volumosas ($N = 1.000$), mesmo pequenas assimetrias são prontamente detectadas pelo teste de Kolmogorov-Smirnov. Esse mesmo comportamento foi observado no único estudo do mapeamento sistemático que aplicou testes estatísticos de hipótese \cite{soares2024}, o qual também rejeitou a normalidade em Go e recorreu a testes não paramétricos.

\subsection{Análise 1 — REST vs. gRPC por Payload (Teste de Mann-Whitney)}

Para responder à \textbf{QP1} (1 fator com 2 tratamentos: REST e gRPC), executou-se o \textbf{Teste de Mann-Whitney} (\textit{Stat $\rightarrow$ Nonparametrics $\rightarrow$ Mann-Whitney}) em cada tamanho de payload ($\alpha = 0{,}05$). A Tabela~5 reúne os resultados e a Figura~6 reproduz integralmente as saídas textuais geradas no Minitab.

\begin{table}[H]
\centering
\caption{Resultados do Teste de Mann-Whitney comparando REST vs. gRPC por tamanho de payload ($\alpha = 0{,}05$)}
\footnotesize
\setlength{\tabcolsep}{3.5pt}
\renewcommand{\arraystretch}{1.35}
\begin{tabularx}{\textwidth}{>{\raggedright\arraybackslash}p{1.15cm}*{6}{>{\raggedleft\arraybackslash}p{1.48cm}}>{\raggedright\arraybackslash}X}
\hline
\rowcolor{cabecalho}
\textbf{Payload} & \textbf{Med. REST (ms)} & \textbf{Med. gRPC (ms)} & \textbf{$\hat{\eta}_{\text{R}}-\hat{\eta}_{\text{G}}$ (ms)} & \textbf{Estat. $W$} & \textbf{$p$-valor} & \textbf{Ganho gRPC} & \textbf{Conclusão} \\
\hline
1 KB   & 0,0405  & 0,0403 & $+0{,}0008$  & 1.033.151,5 & 0,0115       & $+0{,}62\%$  & \textbf{gRPC} (rejeita $H_0$) \\
100 KB & 0,4971  & 0,1958 & $+0{,}2992$  & 1.493.842,0 & $< 0{,}0001$ & $+60{,}60\%$ & \textbf{gRPC} (rejeita $H_0$) \\
1 MB   & 4,6778  & 1,7422 & $+2{,}9320$  & 1.500.500,0 & $< 0{,}0001$ & $+62{,}76\%$ & \textbf{gRPC} (rejeita $H_0$) \\
4 MB   & 18,7834 & 6,9928 & $+11{,}7458$ & 1.499.502,0 & $< 0{,}0001$ & $+62{,}77\%$ & \textbf{gRPC} (rejeita $H_0$) \\
\hline
\end{tabularx}
\fonte{elaborado pelo autor (2026).}
\end{table}

\begin{figure}[H]
\centering
\caption{Saídas do Teste de Mann-Whitney obtidas no Minitab comparando REST e gRPC em cada tamanho de payload}
\begin{minipage}[t]{0.485\textwidth}
  \minitabbox[\linewidth]{%
    \ttfamily\scriptsize
    \textbf{\color{minitabAzul}Mann-Whitney: REST\_1KB; gRPC\_1KB}\\[3pt]
    \begin{tabular}{@{}lrr@{}}
    \textbf{Sample} & \textbf{N} & \textbf{Median} \\
    \hline
    REST\_1KB & 1000 & 0.04050 \\
    gRPC\_1KB & 1000 & 0.04025 \\
    \end{tabular}\\[4pt]
    \textbf{Estimation for Difference ($\eta_1 - \eta_2$)}\\
    Difference: \textbf{0.000793}\\
    95.00\% CI: (0.000167; 0.001417)\\[4pt]
    \textbf{Test} ($H_0: \eta_1-\eta_2=0$ vs $H_1: \eta_1-\eta_2\neq 0$)\\
    W-Value (adj. for ties): \textbf{1033151.50}\\
    P-Value (adj. for ties): \textbf{0.011}
  }
  \par\vspace{3pt}\centering{\footnotesize (a) Payload de 1 KB}
\end{minipage}\hfill
\begin{minipage}[t]{0.485\textwidth}
  \minitabbox[\linewidth]{%
    \ttfamily\scriptsize
    \textbf{\color{minitabAzul}Mann-Whitney: REST\_100KB; gRPC\_100KB}\\[3pt]
    \begin{tabular}{@{}lrr@{}}
    \textbf{Sample} & \textbf{N} & \textbf{Median} \\
    \hline
    REST\_100KB & 1000 & 0.497063 \\
    gRPC\_100KB & 1000 & 0.195834 \\
    \end{tabular}\\[4pt]
    \textbf{Estimation for Difference ($\eta_1 - \eta_2$)}\\
    Difference: \textbf{0.299249}\\
    95.00\% CI: (0.296875; 0.301458)\\[4pt]
    \textbf{Test} ($H_0: \eta_1-\eta_2=0$ vs $H_1: \eta_1-\eta_2\neq 0$)\\
    W-Value (adj. for ties): \textbf{1493842.00}\\
    P-Value (adj. for ties): \textbf{0.000}
  }
  \par\vspace{3pt}\centering{\footnotesize (b) Payload de 100 KB}
\end{minipage}

\vspace{8pt}

\begin{minipage}[t]{0.485\textwidth}
  \minitabbox[\linewidth]{%
    \ttfamily\scriptsize
    \textbf{\color{minitabAzul}Mann-Whitney: REST\_1MB; gRPC\_1MB}\\[3pt]
    \begin{tabular}{@{}lrr@{}}
    \textbf{Sample} & \textbf{N} & \textbf{Median} \\
    \hline
    REST\_1MB & 1000 & 4.67779 \\
    gRPC\_1MB & 1000 & 1.74217 \\
    \end{tabular}\\[4pt]
    \textbf{Estimation for Difference ($\eta_1 - \eta_2$)}\\
    Difference: \textbf{2.93204}\\
    95.00\% CI: (2.91512; 2.94846)\\[4pt]
    \textbf{Test} ($H_0: \eta_1-\eta_2=0$ vs $H_1: \eta_1-\eta_2\neq 0$)\\
    W-Value (adj. for ties): \textbf{1500500.00}\\
    P-Value (adj. for ties): \textbf{0.000}
  }
  \par\vspace{3pt}\centering{\footnotesize (c) Payload de 1 MB}
\end{minipage}\hfill
\begin{minipage}[t]{0.485\textwidth}
  \minitabbox[\linewidth]{%
    \ttfamily\scriptsize
    \textbf{\color{minitabAzul}Mann-Whitney: REST\_4MB; gRPC\_4MB}\\[3pt]
    \begin{tabular}{@{}lrr@{}}
    \textbf{Sample} & \textbf{N} & \textbf{Median} \\
    \hline
    REST\_4MB & 1000 & 18.7834 \\
    gRPC\_4MB & 1000 &  6.9928 \\
    \end{tabular}\\[4pt]
    \textbf{Estimation for Difference ($\eta_1 - \eta_2$)}\\
    Difference: \textbf{11.7458}\\
    95.00\% CI: (11.7189; 11.7725)\\[4pt]
    \textbf{Test} ($H_0: \eta_1-\eta_2=0$ vs $H_1: \eta_1-\eta_2\neq 0$)\\
    W-Value (adj. for ties): \textbf{1499502.00}\\
    P-Value (adj. for ties): \textbf{0.000}
  }
  \par\vspace{3pt}\centering{\footnotesize (d) Payload de 4 MB}
\end{minipage}
\fonte{elaborado pelo autor (2026), transcrito da saída do Minitab (\texttt{GRPCxREST.mpx}).}
\end{figure}

Rejeita-se $H_{0,\text{arq}}$ nos quatro tamanhos ($p < 0{,}05$), mas com comportamento marcadamente distinto entre cargas mínimas e cargas médias/grandes: em \textbf{1~KB}, embora o teste de Mann-Whitney aponte diferença estatisticamente significativa ($p = 0{,}0115$), a diferença prática entre as medianas é residual (0,0405~ms no REST contra 0,0403~ms no gRPC; redução de apenas 0,62\%, $\hat{\eta}_{\text{R}}-\hat{\eta}_{\text{G}} = +0{,}0008$~ms), configurando empate técnico — o custo fixo de enquadramento HTTP/2 e controle de \textit{streams} do gRPC praticamente neutraliza a economia de 333~bytes do Protobuf. A partir de \textbf{100~KB}, o custo de serialização textual JSON domina o tempo de resposta e o gRPC assume ampla liderança, reduzindo a latência mediana em 60,60\% (100~KB), 62,76\% (1~MB, sem qualquer sobreposição entre as amostras: $U = 1.000.000$) e 62,77\% (4~MB).

\subsection{Análise 2 — Efeito do Tamanho do Payload (Kruskal-Wallis e Pós-Teste)}

Para responder à \textbf{QP2} (1 fator com 4 tratamentos: 1~KB, 100~KB, 1~MB e 4~MB), aplicou-se o \textbf{Teste de Kruskal-Wallis} (\textit{Stat $\rightarrow$ Nonparametrics $\rightarrow$ Kruskal-Wallis}) separadamente em cada arquitetura, seguido por comparações par-a-par com \textbf{Mann-Whitney}, conforme a árvore de decisão. Os resultados constam na Tabela~6 e na saída textual do Minitab reproduzida na Figura~7.

\begin{table}[H]
\centering
\caption{Resultados do Teste de Kruskal-Wallis e pós-teste de Mann-Whitney para o fator Tamanho de Payload ($\alpha = 0{,}05$)}
\small
\renewcommand{\arraystretch}{1.35}
\begin{tabularx}{\textwidth}{>{\raggedright\arraybackslash}p{3.0cm}>{\centering\arraybackslash}p{2.2cm}>{\centering\arraybackslash}p{1.3cm}>{\centering\arraybackslash}p{2.3cm}>{\raggedright\arraybackslash}X}
\hline
\rowcolor{cabecalho}
\textbf{Arquitetura} & \textbf{Estatística $H$} & \textbf{GL ($DF$)} & \textbf{$p$-valor} & \textbf{Pós-teste (Mann-Whitney)} \\
\hline
REST (HTTP/1.1) & 3.748,29 & 3 & $< 0{,}0001$ & Todos os 6 pares significativos ($p < 0{,}0001$) \\
gRPC (HTTP/2)   & 3.749,06 & 3 & $< 0{,}0001$ & Todos os 6 pares significativos ($p < 0{,}0001$) \\
\hline
\end{tabularx}
\fonte{elaborado pelo autor (2026).}
\end{table}

\begin{figure}[H]
\centering
\caption{Saídas do Teste de Kruskal-Wallis obtidas no Minitab avaliando o efeito do tamanho do payload em cada arquitetura}
\begin{minipage}[t]{0.485\textwidth}
  \minitabbox[\linewidth]{%
    \ttfamily\scriptsize
    \textbf{\color{minitabAzul}Kruskal-Wallis: Latencia\_ms vs Payload (REST)}\\[3pt]
    \begin{tabular}{@{}lrrrr@{}}
    \textbf{Payload} & \textbf{N} & \textbf{Median} & \textbf{Mean Rank} & \textbf{Z-Value} \\
    \hline
    1KB     & 1000 &  0.0405 &  500.5 & $-47.43$ \\
    100KB   & 1000 &  0.4971 & 1500.5 & $-15.81$ \\
    1MB     & 1000 &  4.6778 & 2501.0 &  $15.83$ \\
    4MB     & 1000 & 18.7834 & 3500.0 &  $47.41$ \\
    Overall & 4000 &         & 2000.5 &          \\
    \end{tabular}\\[4pt]
    \textbf{Test} ($H_0$: All medians are equal)\\
    DF = 3 \quad H-Value = \textbf{3748.29} \quad P-Value = \textbf{0.000}\\
    (Not adjusted and adjusted for ties)
  }
  \par\vspace{3pt}\centering{\footnotesize (a) Arquitetura REST}
\end{minipage}\hfill
\begin{minipage}[t]{0.485\textwidth}
  \minitabbox[\linewidth]{%
    \ttfamily\scriptsize
    \textbf{\color{minitabAzul}Kruskal-Wallis: Latencia\_ms vs Payload (gRPC)}\\[3pt]
    \begin{tabular}{@{}lrrrr@{}}
    \textbf{Payload} & \textbf{N} & \textbf{Median} & \textbf{Mean Rank} & \textbf{Z-Value} \\
    \hline
    1KB     & 1000 & 0.04025 &  500.5 & $-47.43$ \\
    100KB   & 1000 & 0.19583 & 1500.5 & $-15.81$ \\
    1MB     & 1000 & 1.74217 & 2500.5 &  $15.81$ \\
    4MB     & 1000 & 6.99279 & 3500.5 &  $47.43$ \\
    Overall & 4000 &         & 2000.5 &          \\
    \end{tabular}\\[4pt]
    \textbf{Test} ($H_0$: All medians are equal)\\
    DF = 3 \quad H-Value = \textbf{3749.06} \quad P-Value = \textbf{0.000}\\
    (Not adjusted and adjusted for ties)
  }
  \par\vspace{3pt}\centering{\footnotesize (b) Arquitetura gRPC}
\end{minipage}
\fonte{elaborado pelo autor (2026), transcrito da saída do Minitab (\texttt{GRPCxREST.mpx}).}
\end{figure}

Rejeita-se $H_{0,\text{pay}}$ em ambas as arquiteturas ($p < 0{,}0001$), e todos os seis pares de tamanhos diferem significativamente ($p < 0{,}0001$). Além disso, ao passar de 1~KB para 4~MB, a mediana do REST cresceu \textbf{463,8 vezes}, enquanto a do gRPC cresceu \textbf{173,7 vezes}, confirmando escalabilidade substancialmente mais favorável do gRPC para grandes volumes estruturados.

% =============================================================================
% 4. DISCUSSAO E AMEACAS A VALIDADE
% =============================================================================
\section{DISCUSSÃO E AMEAÇAS À VALIDADE}

Os resultados conectam-se diretamente às evidências levantadas no mapeamento sistemático \cite{coimbra2026}. O ganho expressivo do gRPC entre 100~KB e 4~MB (acima de 60\%) converge com \citeonline{buzhin2023}, \citeonline{ain2025} e \citeonline{ye2026}, enquanto o empate técnico observado em 1~KB (diferença prática de apenas 0,0002~ms) dialoga com \citeonline{bolanowski2022}, evidenciando que, em mensagens muito pequenas sem latência de rede e sem TLS, o \textit{overhead} fixo de enquadramento HTTP/2 anula quase integralmente o benefício da serialização binária. Por outro lado, o contraste em relação a \citeonline{soares2024} — que observaram inversão favorável ao REST em 781~KB em Go — reforça que a estrutura do payload (objetos estruturados tipados vs. blocos textuais), o reuso de conexões, a arquitetura de hardware e as versões do runtime influenciam diretamente o comportamento das curvas.

Quanto às ameaças à validade \cite{wohlin2012}, destacam-se: (i) \textbf{Validade de Construto}: comparou-se REST sobre HTTP/1.1 com gRPC sobre HTTP/2, combinando o efeito da serialização ao da versão do protocolo de transporte; (ii) \textbf{Validade Interna}: a execução em sistema operacional multiprogramado sujeita as medições a pausas de coletor de lixo e escalonamento, mitigada por isolamento do ambiente, aquecimento prévio, temporizador monotônico de alta resolução, elevado número de repetições ($N = 1.000$) e estatística não paramétrica; (iii) \textbf{Validade Externa}: o uso de rede local (\textit{loopback}), sem TLS e sem concorrência, limita a generalização direta para redes distribuídas reais; e (iv) \textbf{Validade de Conclusão}: mitigada pelo elevado poder amostral e pela estrita conformidade com a árvore de decisão de \citeonline{sirqueira2015}.

% =============================================================================
% 5. CONCLUSAO
% =============================================================================
\section{CONCLUSÃO}

Este estudo realizou uma análise estatística rigorosa comparando REST (HTTP/1.1 + JSON) e gRPC (HTTP/2 + \textit{Protocol Buffers}) em Go ao longo de quatro tamanhos de payload (1~KB, 100~KB, 1~MB e 4~MB), sobre 8.000 medições coletadas em macOS (Apple M3) e analisadas no Minitab segundo a árvore de decisão de \citeonline{sirqueira2015}. Os testes não paramétricos de Mann-Whitney e Kruskal-Wallis confirmaram ($p < 0{,}05$) que o impacto prático depende fortemente do tamanho da mensagem: em 1~KB, as arquiteturas apresentaram empate técnico (redução mediana de apenas 0,62\% a favor do gRPC); a partir de 100~KB, o gRPC alcançou reduções medianas expressivas de 60,60\% a 62,77\%, além de economizar cerca de 30\% em bytes trafegados e escalar muito melhor até 4~MB. Como trabalhos futuros para o TCC, prevê-se incluir REST sobre HTTP/2, realizar varredura mais fina entre 1~KB e 100~KB e avaliar cenários com latência de rede controlada, TLS e múltiplos níveis de concorrência.

% =============================================================================
% REFERENCIAS
% =============================================================================
\bibliography{referencias}

\end{document}
